% STATUS: conditional framework; outstanding model-specific work is explicit.
\section{P12: Hopf transversality (unfinished application)}
\label{sec:P12}
The source has no derivative calculation. The formulas below specify
the standard reduction; the crossing on P11's path remains to be verified.
\subsection{Equilibrium derivative}\label{slot:P12-1}
For $B(x,g)=a-(k+h)x+(c-g)/(1+x)=0$ with $a,k,c,h$ fixed,
implicit differentiation gives
\[
 B_x=-(k+h)-\frac{c-g}{(1+x)^2},\qquad
 x_g=\frac{1}{(1+x)B_x}\quad(B_x\ne0).
\]
The implicit-function theorem applies locally at P11's simple upper root;
it fails in this parameterization at a scalar fold
\cite[Ch.~1]{GuckenheimerHolmes1983}.
\subsection{Total crossing derivative}\label{slot:P12-2}
Differentiate P04's coefficients along the moving equilibrium:
\[
 \frac{d\lambda_m}{dg}
 =d_g+C_g e^{-\ii\theta_m}+G_g e^{2\ii\theta_m}-H_g e^{-2\ii\theta_m}.
\]
Here every subscript denotes a total derivative, including $x_g$.
\subsection{Outstanding nondegeneracy}\label{slot:P12-3}
Evaluate and certify $\Ree(d\lambda_j/dg)\ne0$ at P11's exact tuple,
with the orientation of increasing $g$ explicit. This is the crossing
hypothesis in \cite{KuznetsovHopf2006}; the spectral conditions alone
do not prove it.
